\chapter{Zielmengen, Einschränkungen und Erweiterungen}
\label{b07guide:ch:restrictions}

\section{Zielmengen anpassen}

\begin{theorem}[Die surjektive Zielmenge ist eindeutig bestimmt]
\label{b07guide:thm:codomain}
Seien $f\colon A\to B$ und $D$ eine Menge mit $f[A]\subseteq D$.
Für denselben Graphen gilt
\[
 f\colon A\twoheadrightarrow D
 \quad\Longleftrightarrow\quad D=f[A].
\]
Insbesondere ist jede Funktion auf ihr Bild surjektiv:
\[
 f\colon A\twoheadrightarrow f[A].
\]
Ist bereits $f\colon A\twoheadrightarrow B$, so ist dieser Graph
auf eine zulässige Zielmenge $D$ genau dann surjektiv, wenn $D=B$.
Eine echte Vergrößerung der Zielmenge erhält Surjektivität nicht.
\end{theorem}

\section{Einschränkungen und das Treffen sämtlicher Fasern}

\begin{theorem}[Genaues Kriterium für surjektive Einschränkungen]
\label{b07guide:thm:restriction}
Seien $f\colon A\to B$, $U\subseteq A$ und $V\subseteq B$.
Dann gilt
\[
 f|_U\colon U\twoheadrightarrow V
 \quad\Longleftrightarrow\quad
 \left\{
 \begin{array}{l}
 U\subseteq f^{-1}[V],\\
 \forall b\in V\quad U\cap\operatorname{Fib}_f(b)\ne\varnothing.
 \end{array}\right.
\]
Die erste Bedingung stellt sicher, daß alle Werte in $V$ liegen.
Insbesondere ist $f|_U\colon U\to B$ genau dann surjektiv,
wenn $U$ jede Faser über $B$ trifft.
\end{theorem}

\begin{theorem}[Einschränkung auf das Urbild einer Zielteilmenge]
\label{b07guide:thm:preimagerestriction}
Seien $f\colon A\twoheadrightarrow B$ und $V\subseteq B$. Dann ist
\[
 f|_{f^{-1}[V]}\colon f^{-1}[V]\twoheadrightarrow V
\]
surjektiv. Dies gilt auch für $V=\varnothing$.
\end{theorem}

\begin{theorem}[Entfernen von Ausgangselementen]
\label{b07guide:thm:deletion}
Seien $f\colon A\twoheadrightarrow B$ und $D\subseteq A$.
Die Einschränkung auf $A\setminus D$ ist genau dann auf $B$
surjektiv, wenn keine Faser ganz in $D$ liegt:
\[
 f|_{A\setminus D}\colon A\setminus D\twoheadrightarrow B
 \quad\Longleftrightarrow\quad
 \forall b\in B\quad \operatorname{Fib}_f(b)\not\subseteq D.
\]
Für ein einzelnes $a\in A$ gilt insbesondere
\[
 f|_{A\setminus\{a\}}\text{ surjektiv auf }B
 \quad\Longleftrightarrow\quad
 \exists x\in A\setminus\{a\}\quad f(x)=f(a).
\]
\end{theorem}

\section{Fortsetzen und zusammenfügen}

\begin{theorem}[Erhalt bei gegebener Fortsetzung]
\label{b07guide:thm:extension}
Seien $A\subseteq C$, $f\colon A\twoheadrightarrow B$ und
$F\colon C\to B$ mit $F|_A=f$. Dann ist $F$ surjektiv auf $B$.

Eine solche Fortsetzung existiert genau dann, wenn
$C=A$ oder $B\ne\varnothing$.
Ist $b_0\in B$ gegeben, liefert die Vorschrift
\[
 F(x)=
 \begin{cases}
 f(x),&x\in A,\\
 b_0,&x\in C\setminus A
 \end{cases}
\]
eine surjektive Fortsetzung.
\end{theorem}

\begin{theorem}[Zusammenfügen einer gegebenen Funktionenfamilie]
\label{b07guide:thm:gluing}
Sei $(f_i\colon A_i\to B)_{i\in I}$ eine gegebene Familie mit
\[
 f_i(x)=f_j(x)\qquad
 (i,j\in I,\ x\in A_i\cap A_j).
\]
Für die dadurch bestimmte Funktion
\[
 F=\bigcup_{i\in I}f_i\colon\bigcup_{i\in I}A_i\to B
\]
gilt
\[
 F\text{ surjektiv auf }B
 \quad\Longleftrightarrow\quad
 \bigcup_{i\in I}f_i[A_i]=B.
\]
Insbesondere genügt ein surjektives $f_i$.
Sind allgemeiner $B_i\subseteq B$ mit
$\bigcup_{i\in I}B_i=B$ und $f_i[A_i]=B_i$ für alle $i$,
so ist $F$ surjektiv. Für $I=\varnothing$ ist $F$ genau dann
surjektiv, wenn $B=\varnothing$.
\end{theorem}